\subsection{Способы находить датасеты (минимум 3 источника)}

Машинное обучение начинается с данных.
Важно, чтобы их было достаточно много и чтобы они были достаточно качественными.
Чем сложнее задача, тем больше данных нужно, чтобы её решить

Бороться с проблемой нехватки данных можно двумя способами.
Первый -- использование краудсорсинга, то есть привлечение людей, готовых разметить много данных.
Во многих ситуациях (например, при оценке поисковой выдачи) без дополнительной разметки никак не обойтись.
Мы расскажем про краудсорсинг подробнее в соответствующем параграфе.
Некоторые проекты, в первую очередь научные и социальные, используют также citizen science -- разметку данных
волонтёрами без какого-либо вознаграждения, просто за чувство причастности к доброму делу исследования.

Второй же способ состоит в использовании неразмеченных данных.
Например, в задаче аннотирования изображений у нас есть огромное количество никак не связанных друг с другом
изображений и текстов.
Однако мы можем использовать их для того, чтобы помочь компьютеру понять, какие слова в принципе могут стоять рядом

\vspace{5mm}
\underline{Публичные репозитории и платформы данных}.
Крупные специализированные платформы, где датасеты выложены для свободного или исследовательского использования
\begin{itemize}
    
    \item \textbf{Kaggle Datasets} содержит огромную коллекция датасетов на самые разные темы (от финансов до
    компьютерного зрения), часто сопровождаемая обсуждениями и примерами анализа (ноутбуками).
    Внутри можно, например, найти набор данных \texttt{titanic} для прогнозирования выживания пассажиров.
    
    \item \textbf{UCI Machine Learning Repository} есть один из старейших и самых уважаемых источников классических
    датасетов, широко используемых в академических исследованиях и для обучения.
    Например, датасет \texttt{Iris} для задач классификации.

\end{itemize}

\underline{API и порталы государственных/международных организаций}.
Многие государственные и международные структуры публикуют данные в открытом доступе.
\begin{itemize}
    
    \item \textbf{Порталы открытых данных}, тот же \texttt{data.gov.ru}, содержат огромные массивы данных от
    государственных ведомств: демография, экономика, экология, транспорт.
    Например, данные о государственных закупках или о качестве атмосферного воздуха в городе
    
    \item \textbf{Всемирный банк, ВОЗ, Eurostat} предоставляют статистику по странам мира в разных областях.
    Данные часто доступны для скачивания в структурированном виде (CSV, JSON).
    Например, статистика ВВП по всем странам за последние 60 лет

\end{itemize}

\underline{Специализированные поисковые системы и агрегаторы}
Инструменты, созданные специально для поиска по публичным датасетам.
\begin{itemize}
    
    \item \textbf{Google Dataset Search} -- поисковая система, которая индексирует датасеты с различных источников.
    Работает по аналогии с обычным поиском, но ищет именно метаданные датасетов
    
    \item \textbf{Awesome Public Datasets (на GitHub)} -- это кураторские списки (репозитории), где датасеты
    сгруппированы по темам (биология, география, социальные науки и т.д.).

\end{itemize}